\section{Measuring matchertext containment}
\label{sec:prevent}

\Cref{sec:resist} argued that matchertext \emph{can} contain injection; this
section asks \emph{how much} of real injection it reaches, with two measurements
over the corpus of \cref{sec:corpus}: a simulation replaying every recorded payload
against a matchertext-aware host language, and one executed such host
language: a modified SQLite, running the SQL payloads through a live parser
(\cref{sec:prevent:exec}).

Every payload falls into one of three outcomes, and matchertext helps in two by
different routes. Under \emph{inert embedding} the value is valid matchertext and
sits as data inside its matcher-delimited slot, metacharacters read as text; under
\emph{rejection} it carries an unmatched matcher, fails the passive \textsc{Verify}
scan, and never embeds. In the third, \emph{outside}, matchertext cannot help: the
value lands in a context no matcher pair delimits, or a sink the host language runs
by design. Which outcome an attack falls into is a property of its payload, and so
measurable.

\begin{figure*}[t]
\footnotesize
%% GENERATED by pipeline/latex.py -- do not edit.
%% Regenerate with: python3 pipeline/latex.py
\begin{verbatim}
                      11,265 payload-backed CVEs
                    ┌──────────────┴──────────────┐
                   yes                            no
       matcher-delimitable context?            OUTSIDE
               9,275 (82%)                     1,990 (18%)
        ┌───────────┴───────────┐        ┌──────────┴──────────┐
      pass                    fail       807                  1183
    VERIFY(v)               VERIFY(v)  execution sinks       non-matcher delimiters
        ▼                       ▼      templates, eval,      shell ; | &
   INERT EMBED              REJECTED   javascript:           CR/LF, XML <>
   8,761 (94%)              514 (6%)      (A4)
\end{verbatim}

\caption{Containment outcome as two decisions. The context condition is tested
first, so the \dataOutsideCVEs outside are never candidates rather than escapes; the value
condition then separates inert embedding from rejection. The leaves reconcile
with \cref{tab:mech}.}
\label{fig:prevent:tree}
\end{figure*}

\paragraph{Simulating containment.}
For each of the \dataPayloadBackedCVEs payload-backed CVEs, we reuse its syntactic
skeleton and decide whether the target syntax is matcher-delimitable and, if so,
whether the skeleton passes \textsc{Verify}, contained inertly if it does and by
rejection if it does not. This simulates each payload against a matchertext-aware
host language, running \textsc{Verify} over the skeleton without sending it through
a live parser.

\begin{table}[t]
\centering
\small
%% GENERATED by pipeline/latex.py -- do not edit.
%% Regenerate with: python3 pipeline/latex.py
\begin{tabular}{lrrr}
\toprule
Syntax & Inert embed & Rejected & Outside \\
\midrule
HTML / DOM (XSS)    & 4{,}424 &  34 & 270 \\
SQL                 & 4{,}197 & 477 &   0 \\
LDAP                &       0 &   1 &   0 \\
NoSQL               &       1 &   0 &   0 \\
XPath / XQuery      &       1 &   0 &   0 \\
\midrule
Shell command       & --- & --- & 1{,}086 \\
Code / eval         & --- & --- & 441 \\
CRLF / header       & --- & --- &  52 \\
Formula (CSV)       & --- & --- &  37 \\
Template            & --- & --- &  29 \\
Argument            & --- & --- &  15 \\
Expression language & --- & --- &  12 \\
XML                 & --- & --- &  10 \\
\midrule
Total & 8{,}623 & 512 & 1{,}952 \\
\bottomrule
\end{tabular}

\caption{Containment outcome by syntax over all \dataPayloadBackedCVEs
payload-backed CVEs. Delimitable contexts (upper block) are contained by inert
embedding or rejection; \emph{Outside} counts CVEs failing a premise of
\cref{sec:threat}.}
\label{tab:mech}
\end{table}

\paragraph{Results.}
Of the \dataPayloadBackedCVEs payload-backed injection CVEs, the recorded
payloads of \textbf{\dataContained} (\textbf{\dataContainedShare}) are contained
by a matcher-delimited slot that a matchertext-aware host language reads to
matcher balance, assuming no intervening decoding.
The split turns on the host language's native delimiter, and the per-syntax
mechanisms of \cref{sec:resist} appear directly in the corpus. LDAP is the pure
rejection case: every real LDAP payload fails \textsc{Verify} (\dataLdapInert
inert, \dataLdapRejected rejected), its parenthesis breakout carrying an unmatched
\verb|)|. SQL and XSS break out on nonmatchers, so they are overwhelmingly valid
matchertext and embed inertly, the \dataSqlRejected SQL rejections carrying an
unbalanced function or subquery parenthesis.
The ratio is descriptive, not predictive: these payloads were written against
quote- and tag-delimited host languages, and an attacker aware of a
matcher-delimited slot would balance the matchers, shifting cases from rejection to
inert embedding without affecting containment, since both routes hold the value.

\paragraph{Where matchertext cannot help.}
The remaining \dataOutsideCVEs payload-backed CVEs are outside both conditions,
split by the premise of \cref{sec:threat} they fail. \dataOutsideSemantic target
execution/semantic sinks (server-side templates, expression languages, code
evaluation, spreadsheet formulas, \verb|javascript:| URLs) that the host language
runs by design: the template payload \verb|{{7*7}}| is already valid matchertext
yet still executes. The other \dataOutsideNonmatcher target non-matcher contexts no
matcher pair bounds: shell separators (\verb|;|~\verb-|-~\verb|&|), CR/LF header
breaks, and XML/XXE angle brackets.

\paragraph{From the measured subset to the whole corpus.}
The outcomes so far cover only the \dataPayloadBackedCVEs payload-backed CVEs.
How far does the reach extend across all \dataInjectionCVEs injection CVEs, most
carrying no payload to replay? We estimate it by \emph{projection}: carrying the
per-syntax containment rates from the payload-backed subset onto the full corpus,
weighted by each syntax's CVE count.
Classifying the \emph{mechanism} needs a payload, but the \emph{context}
condition, whether a syntax is matcher-delimitable at all, reads from the syntax
label alone, so it projects directly.
At most \textbf{\dataHostableCVEs} of the \textbf{\dataInjectionCVEs} injection
CVEs (\textbf{\dataHostableShare}) target a matcher-delimitable syntax.
This is an upper bound, since some HTML/DOM cases are \verb|javascript:| semantic
sinks represent \textbf{\dataJsSinkRate} (\dataHtmlDomJsSinks) of payloads.

Projecting the rate needs care, because payload-backed CVEs are not a random
sample. A SQL payload is a literal string a write-up can quote, whereas code-eval
and argument payloads are built at run time and rarely appear verbatim, so the
syntaxes matchertext can delimit are over-represented among CVEs that carry a
payload. Naively applying the pooled contained rate of \dataContainedShare to the whole corpus would
therefore overcount, yielding \dataNaiveEstimate contained CVEs
(\dataSkewCost above the stratified result).
To correct for this bias we post-stratify: compute the contained rate
separately within each syntax, then weight each rate by that syntax's
corpus share (\cref{tab:strata}).
The resulting central estimate is
\textbf{\dataCentralEstimate} (\textbf{\dataCentralShare}), below the
\dataHostableShare ceiling.

\begin{table}[t]
\centering
\small
%% GENERATED by pipeline/latex.py -- do not edit.
%% Regenerate with: python3 pipeline/latex.py
\begin{tabular}{lrrr}
\toprule
Syntax & CVEs & Payload-backed & Contained \\
\midrule
HTML / DOM (XSS)    &  50{,}884 & 5{,}828 & 94.0\% \\
SQL                 &  22{,}225 & 5{,}157 & 100.0\% \\
Shell command       &   9{,}672 & 1{,}696 & 0.0\% \\
Code / eval         &   6{,}666 &     626 & 0.0\% \\
CRLF / header       &       656 &      83 & 0.0\% \\
Argument            &       430 &      27 & 0.0\% \\
Formula (CSV)       &       321 &      48 & 0.0\% \\
Expression language &       273 &      20 & 0.0\% \\
\bottomrule
\end{tabular}

\caption{Post-stratification strata: per-syntax rates from the payload-backed
subset, weighted by each syntax's corpus share. A dagger would mark a syntax
with no payload-backed CVE, falling back to the structural rule of
\cref{sec:threat}.}
\label{tab:strata}
\end{table}
Against a single denominator, the full \dataInjectionCVEs injection CVEs, the reach
runs from a floor of \textbf{\dataFloorShare} to a ceiling of
\textbf{\dataHostableShare}, with \dataCentralShare central. The floor counts only
the \dataContained CVEs with a payload we verified as contained and treats every
CVE without a recoverable payload as uncontained, a hard lower bound; the
\dataContainedShare above is instead the contained share \emph{within} the
payload-backed subset.

\paragraph{Does the payload-backed subset cover enough to trust?}
Only \dataPayloadBackedCVEs of the \dataInjectionCVEs injection CVEs carry a
replayable payload. That gap is
payload availability, a property of the disclosure ecosystem, not a containment
failure: Mandiant finds \textbf{51\%} of 2021--2022 vulnerabilities eventually
acquired public exploit or PoC code~\cite{mandiant23tte} and VulnCheck reports
\textbf{26\%} for 2025 identifiers~\cite{vulncheck26exploit}, while our corpus is
stricter still, \dataInjectionWithPoc (\dataInjectionWithPocShare) carrying any PoC
reference at all. The measurable population is therefore roughly half the corpus,
and the fragment-level bound should be read against it, not against
\dataInjectionCVEs.

\subsection{A matchertext-aware SQLite}
\label{sec:prevent:exec}

The measurement so far is a simulation, which leaves the sharpest question open:
can a \emph{real} parser and interpreter hold a payload the way the discipline
promises?
\Cref{sec:threat} isolates the two premises a simulation cannot reach, parse
inertness (L2) and interpreter inertness (L3), each a property of a concrete
implementation.
To test them we build one matchertext-aware host language, a modified SQLite
engine,\footnote{The fork and the driver of this experiment are available in an
anonymized repository:
\url{https://anonymous.4open.science/r/matchertext-sqlite-4301/}.} and run the
recorded SQL payloads through it, from the CVE-derived SQL subset and the
standalone SQL corpora.

\paragraph{The modified engine.}
Our driver realizes the discipline of \cref{sec:resist} directly: it places a
value in the SQL text inside an \verb|M'(...)'| literal, wraps an identifier in
\verb|[|$\dots$\verb|]| that the engine reads to matcher balance (A1), and
prepares the statement through the ordinary interface.
The fork offers this hole beside legacy composition rather than replacing it, so
existing SQL keeps working and adoption is incremental. That additive design is the
point-of-use character of A2: a value still concatenated into surrounding SQL,
outside a hole, is unprotected. An opt-in strict mode refuses the legacy prepare
entry points, turning A2 into an invariant the engine enforces; the measurement
reports the additive default and strict build side by side.

\paragraph{Fixtures and oracle.}
We place each of the \dataSqliPayloads distinct SQL payloads (\dataSqliFamilies
families) through every compatible hole-arm pair and compare it against the same
hole holding a benign value.
The 16 fixtures cover single- and double-quoted text, raw and parenthesized
numbers, \texttt{IN} lists, \texttt{ORDER BY} and \texttt{LIMIT} expressions,
one- to four-column selects, one- and two-column inserts, and identifier
positions.
Fixtures and arms are each typed by hole kind, so \dataSqliValueHosts value
fixtures run against \dataSqliValueArms value arms and \dataSqliIdentHosts
identifier fixtures against \dataSqliIdentArms identifier arms:
\dataSqliCasesPerPayload cases per payload, \dataSqliCases in all, the two typed
strict entry points swept separately (\cref{tab:sqli}, right).
The \dataSqliIdentMtMalformed name-hole cases that do not compile are the
\dataSqliNulPayloads payloads with a trailing zero byte, which ends a SQLite
statement wherever it appears rather than the hole failing to contain it. The set
retains every recorded payload.
SQLite variants allow one transport percent-decoding pass, SQL Server
\verb|N'...'| text, PostgreSQL dollar-quoted text, and a terminal MySQL \verb|#|
comment, with no function, operator, query, or objective changed.
The oracle is the compiled statement, not the string: the \textsc{Explain}
bytecode with the value operand redacted, the read-only flag, whether any SQL
follows the first statement, and the rows returned, each read against the benign
baseline. A payload counts against a defense only where the control, naive
concatenation into a value context or a bare name, actually breaks out; most
recorded payloads are fragments that never parse against SQLite. That effective set
is \dataSqliEffective payloads.

\begin{table*}[t]
\centering
\small
%% GENERATED by pipeline/latex.py -- do not edit.
%% Regenerate with: python3 pipeline/latex.py
\begin{tabular}{l rrrr rrrr}
\toprule
& \multicolumn{4}{c}{Additive} & \multicolumn{4}{c}{Strict} \\
\cmidrule(lr){2-5} \cmidrule(lr){6-9}
Delivery & Inert & Broke out & Refused & Not SQL & Inert & Broke out & Refused & Not SQL \\
\midrule
Concatenated into host SQL         & 35{,}740 & 2{,}625 & 0 & 53{,}713 & 0 & 0 & 92{,}078 & 0 \\
Concatenated as a name             & 1{,}221 & 21 & 0 & 11{,}912 & 0 & 0 & 13{,}154 & 0 \\
\midrule
\texttt{\%Q} escaping              & 92{,}078 & 0 & 0 & 0 & 0 & 0 & 92{,}078 & 0 \\
\texttt{\%q} escaping              & 92{,}078 & 0 & 0 & 0 & 0 & 0 & 92{,}078 & 0 \\
Bound parameter                    & 92{,}078 & 0 & 0 & 0 & 0 & 0 & 92{,}078 & 0 \\
\midrule
\texttt{M'(\ldots)'} value hole    & 92{,}078 & 0 & 0 & 0 & 76{,}734 & 0 & 15{,}344 & 0 \\
\texttt{M'(\%M)'} verified         & 76{,}734 & 0 & 15{,}344 & 0 & 76{,}734 & 0 & 15{,}344 & 0 \\
\texttt{[\ldots]} name hole        & 13{,}150 & 0 & 0 & 4 & 10{,}958 & 0 & 2{,}192 & 4 \\
\bottomrule
\end{tabular}

\caption{Executed outcome by delivery path over the \dataSqliPayloads SQL
payloads (\dataSqliCases cases), the additive default beside the opt-in strict
build. Matcher-delimited holes (lower block) contain every payload; the
concatenation controls (upper block) are the paths that break out. \emph{Not
SQL} counts fragments that do not parse against SQLite; \emph{Refused},
statements rejected before the parser.}
\label{tab:sqli}
\end{table*}

\paragraph{Results.}
Over the effective set no matchertext hole admitted a breakout (\cref{tab:sqli}):
\dataSqliMtBreakout through the \verb|M'(...)'| value hole, \dataSqliMtStrictBreakout
through the verify-then-embed route, and \dataSqliIdentMtBreakout through the
\verb|[|$\dots$\verb|]| name hole.
The legacy value defenses hold as well, \texttt{\%Q} and \texttt{\%q} escaping
and a bound parameter, which is the baseline.
Concatenation is the contrast: \dataSqliConcatBreakout breakouts of the value
control and \dataSqliIdentConcatBreakout of the identifier control, on the very
same payloads.
That breakout is available only because the additive build leaves legacy
composition open. Recompiled in strict mode, every text-splicing prepare, the
concatenation and the in-band \verb|M'(...)'| alike, is refused before the parser
(\cref{tab:sqli}, right), so the \dataSqliEffective effective-set attacks have no
path at all and A2 holds by construction.
The only admitted paths are then the typed \texttt{?V} value and \texttt{?I}
identifier parameters of \texttt{sqlite3\_matchertext\_prepare\_v3()}, which
contain the corpus as the additive holes do: \dataSqliStrictValueBreakout value and
\dataSqliStrictIdentBreakout identifier breakouts. The \texttt{?V} route binds each
valid value and refuses the \dataSqliStrictValueRejected that are not matchertext;
\texttt{?I} holds each as a single name.
Short of strict mode, an application closes the gap itself with the discipline's
two \textsc{Verify} checks, valid input and a valid assembled statement, which are
exactly the value-hole and name-hole rows; the concatenation control is their
absence. The value hole embeds each payload inertly, the verified route refuses the
\dataSqliMtStrictRejected non-matchertext (inert embedding and rejection,
\cref{sec:resist}), and the identifier hole holds each as a single name, the one
position a bound parameter cannot fill.

\paragraph{Falsifiability.}
A zero is evidence only if the oracle could have reported a one.
We rebuild the scanner deliberately broken, once so the boundary locator stops
at the first closer and once so \textsc{Verify} always accepts, and rerun.
Each broken build changes the verdict of the holes that rest on it, in up to
\dataSqliSabotageDetect of cases. Those are the value and name holes that end by
matcher balance and the verified route that calls \textsc{Verify}. So the honest
zero is read against demonstrated detection power rather than a blind harness.
The C scanner and an independent implementation agree on \textsc{Verify} for
every payload, and the run is deterministic.
In sum, the simulated SQL containment of \cref{tab:mech} survives a live parser and
strict mode makes A2 an invariant the engine enforces; the experiment discharges L2
and L3 for SQL alone, and each further host language needs its own executed check.
